\documentclass[10pt,a4paper]{amsart}
\usepackage[margin=19mm]{geometry}
\usepackage{amsmath,amssymb,mathtools}
\usepackage[scr=rsfs,cal=euler]{mathalfa}
\usepackage{xfrac}
\usepackage[unicode,hidelinks,bookmarks=false]{hyperref}
\newcommand{\ie}{i.e.\ }
\newcommand{\cf}{cf.\ }
\newcommand{\resp}{respectively\ }
\newcommand{\Subsection}{\subsection}
\providecommand{\nobreakdash}{\nobreak\hbox{-}}
\setlength{\emergencystretch}{2em}
\multlinegap0pt
\usepackage{xcolor}
\usepackage{amssymb}
\usepackage{tikz}
\newtheorem{proposition}{Proposition}
\newcommand{\FF}{\mathbb{F}}
\title{The Geometry of Pattern Groups}
\author{Jo\~{a}o Dias, Claudio Alexandre Piedade}
\date{Source: arXiv:2608.11103v1 (CC BY 4.0)}
\begin{document}
\maketitle

\noindent\emph{Source and licence.} The Geometry of Pattern Groups. Version arXiv:2608.11103v1 (CC BY 4.0), distributed by arXiv under the Creative Commons Attribution 4.0 International licence, http://creativecommons.org/licenses/by/4.0/, as recorded in the arXiv licence metadata for this exact version and shown on the abstract page https://arxiv.org/abs/2608.11103v1. Full licence notice: \texttt{licenses/arxiv-2608.11103-CC-BY-4.0.txt}.\par\smallskip

\noindent\textit{Editorial scope.} Counted unit: the article's proposition prop:diag\_props (source0.tex lines 907--914). The article's complete original proof of it is delivered with it (source0.tex lines 915--1009, one \texttt{proof} environment of 95 lines). No mathematics is written, repaired or re-derived: the statement and the proof are the article's own lines, in the article's own notation, and no retained line is substituted or re-flowed.  Retained uncounted as the source-local context the proof needs: lines 152--163; lines 851--853.  Nothing was omitted from any retained span, so the package carries no omission record.  No retained line contains a citation, so the wrapper carries no bibliography and no bibliography entry.  No retained line contains a figure environment, an image-inclusion command or drawing code, so no image is delivered and the package declares no asset dependency.  Editorial formatting only, in addition: an AMS \texttt{amsart} preamble that restates the article's own declarations and macros used by the retained lines, namely the environments \texttt{proposition} on separate counters, together with the standard packages those lines need.  The retained environments print in this document as Proposition 1 in order of appearance, whereas the article prints the same items with the numbering of its own full document.  The complete native source of the article as distributed in the arXiv e-print for this exact version is bundled unchanged as \texttt{sources/207387/source0.tex}.
\begin{proposition}\label{prop:gens_props}
    Let $(i,j)$ and $(k,l)\in\Delta_n$, and $t,t'\in \FF_q\setminus\{0\}$ then:
    \begin{enumerate}
        \item $e_{i,j}(t)e_{k,l}(t')=e_{k,l}(t')e_{i,j}(t)$ if and only if $i\neq l$ and $j\neq k$;
        \item $e_{i,j}(t)e_{i,j}(t')=e_{i,j}(t')e_{i,j}(t)=e_{i,j}(t+t')$;
       
        
        \item let $(i,j),(j,k),(i,k)\in \Delta_n$, then $e_{i,k}(tt')=e_{i,j}(t)e_{j,k}(t')e_{i,j}(-t)e_{j,k}(-t')=[e_{i,j}(t),e_{j,k}(t')]$;
        \item $e_{i,j}^{-1}(t) = e_{i,j}(-t)$
        \item $o(e_{i,j}(t)) = p$;
    \end{enumerate}
\end{proposition}

\begin{proposition}\label{prop:UniT_IP}
Let $K,L\subseteq I$. Then $$\langle e_{i}(t)\mid t\in T, i\in K \rangle \cap \langle e_{i}(t)\mid t\in T, i\in L \rangle = \langle e_{i}(t)\mid t\in T, i\in K\cap L \rangle.$$
\end{proposition}

\begin{proposition}\label{prop:diag_props}
    For $i\in I$ we have that:
    \begin{enumerate}
        \item $ U^{[i,i+1]}_4 = U^{[i+1,i]}_4$;
        \item $U^{[i,i+1]}_3 \neq U^{[i,i+1]}_4$ and $U^{[i+1,i]}_3 \neq U^{[i+1,i]}_4$.
        \item Let $f$ be the minimal positive integer such that $U^{[i,i+1]}_f\cap U^{[i+1,i]}_f\neq U^{[i,i+1]}_{f-1}\cup U^{[i+1,i]}_{f-1}$; if $q\neq 2$ then $f=3$, otherwise, $f=4$.
    \end{enumerate}
\end{proposition}

\begin{proof}
    $\textbf{(1)}$ Let $x\in U^{[i+1,i]}_4 = U^{i+1}U^iU^{i+1}U^i$.
    We have that $U^j=\langle e_j(t)\mid t\in T\rangle$, for any $j\in I$. This means that any element of $U^j$ can be written as $e_j(\alpha)$, with $\alpha\in \FF_q$.
    First, note that if $x$ is in either of the following sets 
    \begin{align*}
        \{&e_i(\alpha_1),\ e_{i+1}(\alpha_1)\mid \alpha_1\in \FF_q\}\\
        \{&e_{i}(\alpha_1)e_{i+1}(\alpha_2),\ e_{i+1}(\alpha_1)e_i(\alpha_2) \mid \alpha_1,\alpha_2\in\FF_q\}\\
        \{&e_{i}(\alpha_1)e_{i+1}(\alpha_2)e_{i}(\alpha_3),\ e_{i+1}(\alpha_1)e_{i}(\alpha_2)e_{i+1}(\alpha_3)\mid \alpha_1,\alpha_2,\alpha_3\in \FF_q\}
    \end{align*}
    we can easily check that $x\in U_4^{[i,i+1]}$.
     Hence, consider then $x=e_{i+1}(\alpha_1)e_{i}(\alpha_2)e_{i+1}(\alpha_3)e_i(\alpha_4)$, where $\alpha_k\in \FF_q\setminus\{0\}$ for all $k=1,2,3,4$.
    By writting an element $e_{j}(\alpha)$ as $(\mathbf{1} + f_{j,j+1}(\alpha))$, we have that $e_{i+1}(\alpha_1)e_{i}(\alpha_2)=\mathbf{1}+f_{i,i+1}(\alpha_2)+f_{i+1,i+2}(\alpha_1)$ and $e_{i+1}(\alpha_3)e_{i}(\alpha_4)=\mathbf{1}+f_{i,i+1}(\alpha_4)+f_{i+1,i+2}(\alpha_3)$. First, we will assume that $\alpha_1+\alpha_3\neq 0$. Therefore, we have that 
    \begin{align*}
        e_{i+1}(\alpha_1)e_{i}(\alpha_2)e_{i+1}(\alpha_3)e_i(\alpha_4) &= (\mathbf{1}+f_{i,i+1}(\alpha_2)+f_{i+1,i+2}(\alpha_1))(\mathbf{1}+f_{i,i+1}(\alpha_4)+f_{i+1,i+2}(\alpha_3))\\
        &= \mathbf{1} + f_{i,i+1}(\alpha_2+\alpha_4) + f_{i+1,i+2}(\alpha_1+\alpha_3) + f_{i,i+2}(\alpha_2\alpha_3)\\
        & = e_{i}(\alpha_2+\alpha_4)e_{i+1}(\alpha_1+\alpha_3)(\mathbf{1} - f_{i,i+2}(A))  
    \end{align*}
    with $A = \alpha_2\alpha_1 + \alpha_4\alpha_1+\alpha_4\alpha_3$.
    Then, we have that 
    \begin{align*}
        e_{i,i+2}(-A) = \mathbf{1} - f_{i,i+2}(A) &= e_{i+1}(-\alpha_1-\alpha_3)e_i(-A(\alpha_1+\alpha_3)^{-1}) e_{i+1}(\alpha_1+\alpha_3)e_i(A(\alpha_1+\alpha_3)^{-1})
    \end{align*}
    Hence, as $e_{i+1}(-\alpha_1-\alpha_3) = e_{i+1}(\alpha_1+\alpha_3)^{-1}$, we have
\begin{align*}
        x&= e_{i}(\alpha_2+\alpha_4)e_{i+1}(\alpha_1+\alpha_3)(\mathbf{1} - f_{i,i+2}(A))\\
        &= e_{i}(\alpha_2+\alpha_4)e_{i+1}(\alpha_1+\alpha_3)e_{i+1}(-\alpha_1-\alpha_3)e_i(-A(\alpha_1+\alpha_3)^{-1}) e_{i+1}(\alpha_1+\alpha_3)e_i(A(\alpha_1+\alpha_3)^{-1})\\
        &= e_{i}(\alpha_2+\alpha_4-A(\alpha_1+\alpha_3)^{-1}) e_{i+1}(\alpha_1+\alpha_3)e_i(A(\alpha_1+\alpha_3)^{-1})
    \end{align*}
which belongs to $U_4^{[i,i+1]}$.
Assume now that $\alpha_1+\alpha_3=0$. In this case we have that 
\begin{align*}
e_{i+1}(\alpha_1)e_{i}(\alpha_2)e_{i+1}(\alpha_3)e_i(\alpha_4) &= (\mathbf{1}+f_i(\alpha_2)+f_{i+1}(\alpha_1))(\mathbf{1}+f_i(\alpha_4)+f_{i+1}(\alpha_3))\\
        &= \mathbf{1} + f_i(\alpha_2+\alpha_4) + f_{i,i+2}(\alpha_2\alpha_3)\\
        & = e_{i}(\alpha_2+\alpha_4)(\mathbf{1} + f_{i,i+2}(\alpha_2\alpha_3)).  
    \end{align*}
    As before, note that we can rewrite $e_{i,i+2}(\alpha_2\alpha_3)= \mathbf{1} + f_{i,i+2}(\alpha_2\alpha_3)$ as the following product $$e_{i}(1)e_{i+1}(\alpha_2\alpha_3)e_i(-1)e_{i+1}(-\alpha_2\alpha_3).$$
    Therefore, we have
\begin{align*}
        x&= e_{i}(\alpha_2+\alpha_4)(\mathbf{1} + f_{i,i+2}(\alpha_2\alpha_3))\\
        &= e_{i}(\alpha_2+\alpha_4)e_{i}(1)e_{i+1}(\alpha_2\alpha_3)e_i(-1)e_{i+1}(-\alpha_2\alpha_3)\\
        &= e_{i}(\alpha_2+\alpha_4+1)e_{i+1}(\alpha_2\alpha_3)e_i(-1)e_{i+1}(-\alpha_2\alpha_3)
    \end{align*}
which also belongs to $U_4^{[i,i+1]}$.
With this, we have proven that $U^{[i+1,i]}_4\subseteq U^{[i,i+1]}_4$
We can use equivalent arguments to prove that $U^{[i,i+1]}_4\subseteq U^{[i+1,i]}_4$.
    

    $\textbf{(2)}$ Now, clearly $U_3^{[i,i+1]}\subset U_4^{[i,i+1]}$. Consider $x\in U_4^{[i,i+1]}$ such that $x=e_{i,i+2}(a)$ and $a\in \FF_q\setminus \{ 0\}$. Indeed, we know that $x\in U_4^{[i,i+1]}$ by Proposition~\ref{prop:gens_props}(3).
    Suppose we could write $x$ as a word in $U_3^{[i,i+1]}$, i.e. $x = e_i(\alpha_1)e_{i+1}(\alpha_2)e_i(\alpha_3)$.
    Therefore, we can express $x$ as 
    \begin{align*}
        x&= \mathbf{1} + f_i(\alpha_1+\alpha_3)+f_{i+1}(\alpha_2)+f_{i,i+2}(\alpha_1\alpha_2).
    \end{align*}
    In order for $x=e_{i,i+2}(a)$, we need $\alpha_1\alpha_2=a$, $\alpha_1+\alpha_3=0$ and $\alpha_2=0$, a contradiction.
    Hence $x\notin U_3^{[i,i+1]}$.
    A similar argument proves the case $U_3^{[i+1,i]}\neq U_4^{[i+1,i]}$

    $\textbf{(3)}$ 
    First, note that the reverse inclusion is always true, as  $U_{f-1}^{[i,i+1]}$ and $U_{f-1}^{[i+1,i]}$ are clearly subsets of both $U_{f}^{[i,i+1]}$ and $U_{f}^{[i+1,i]}$. 
    
    We will first prove that for any pair $\{i,i+1\}$ ($i=1,\ldots,n-1$) and $f=1,2$, we have $U^{[i,i+1]}_f\cap U^{[i+1,i]}_f = U^{[i,i+1]}_{f-1}\cup U^{[i+1,i]}_{f-1}$.
    For $f=1$, we have that $U^{[i,i+1]}_1\cap U_1^{[i+1,i]} = U^i\cap U^{i+1} = \langle \mathbf{1}\rangle$ by Proposition~\ref{prop:UniT_IP}. As, by definition, $U_0^{[k,l]}=B$ and $B=\langle \mathbf{1}\rangle$, we have $U_0^{[i,i+1]}=\langle \mathbf{1}\rangle = U_0^{[i+1,i]}$, giving the equality for $f=1$.
    
    Consider now $f=2$.
    We need to prove that $U^iU^{i+1}\cap U^{i+1}U^{i} \subseteq   U^i \cup U^{i+1}$.
    Consider $x\in U^iU^{i+1}\cap U^{i+1}U^i$.
    Hence, $x$ can be written as both a word $e_i(a)e_{i+1}(b)$ and a word $e_{i+1}(c)e_{i+1}(d)$, i.e. we have that $$\mathbf{1}+f_i(a)+f_{i+1}(b)+f_{i,i+2}(ab) = \mathbf{1} + f_i(d)+f_{i+1}(c).$$
    This is only true when $a=d$, $b=c$ and $ab=0\Leftrightarrow a=0\vee b=0$.
    When $a=0$, we get that $x=e_{i+1}(b)=e_{i+1}(c)\in U^{i+1}$; when $b=0$, we get that $x=e_i(a)=e_i(d)\in U^i$. In either case, we have that $x\in U^i\cup U^{i+1}$, as wanted.

    Now, we will consider the case $f=3$. Let $x\in U^iU^{i+1}U^i\cap U^{i+1}U^{i}U^{i+1}$, i.e.
    $$x=e_i(\alpha_1)e_{i+1}(\alpha_2)e_{i}(\alpha_3) = e_{i+1}(\alpha_4)e_{i}(\alpha_5)e_{i+1}(\alpha_6)$$ with $\alpha_k\in \FF_q$ for $k=1,\ldots,6$. 
    Hence, we have 
    \begin{align*}
        e_i(\alpha_1)e_{i+1}(\alpha_2)e_{i}(\alpha_3) &= e_{i+1}(\alpha_4)e_{i}(\alpha_5)e_{i+1}(\alpha_6)\\
        \mathbf{1} + f_i(\alpha_1+\alpha_3)+f_{i+1}(\alpha_2)+f_{i,i+2}(\alpha_1\alpha_2)&=\mathbf{1}+f_i(\alpha_5)+f_{i+1}(\alpha_4+\alpha_6)+f_{i,i+2}(\alpha_5\alpha_6),
    \end{align*}
    which implies that $\alpha_1+\alpha_3=\alpha_5$, $\alpha_2=\alpha_4+\alpha_6$ and $\alpha_1\alpha_2=\alpha_5\alpha_6$.
    Note that if, for some $k$, $\alpha_k=0$, we have trivially that $x\in U^iU^{i+1}\cup U^{i+1}U^i$. So, we will suppose that for all $k=1,\ldots,6$, $\alpha_k\neq 0$.
    
    \textit{Case $q\neq 2$:} First, note that the intersection $U^iU^{i+1}U^i\cap U^{i+1}U^{i}U^{i+1}$ has non-trivial elements which can be written will all $\alpha_k\neq 0$. Indeed, we can take $x$ in this intersection where $\alpha_5=\alpha_2$, which will imply that $\alpha_1=\alpha_6$ and $\alpha_3=\alpha_4$ (all taken to be non-zero).
    Additionally, this means that the matrix of $x$ has a non-null $(i,i+2)$-entry. We remind the reader that, for $\alpha,\beta\in\FF_q$, $e_{i}(\alpha)e_{i+1}(\beta) = \mathbf{1} + f_i(\alpha)+f_{i+1}(\beta)+f_{i,i+2}(\alpha\beta)$, while $e_{i+1}(\beta)e_{i+1}(\alpha) = \mathbf{1} + f_i(\alpha)+f_{i+1}(\beta)$. Since
    any element of $U^{i+1}U^i$ is always a matrix with a null $(i,i+2)$-entry, $x\notin U^{i+1}U^i$. Moreover, given an element from $U^{i}U^{i+1}$, its $(i,i+2)$-entry is the product of its $(i,i+1)$ and $(i+1,i+2)$ entries. However, that would mean $(\alpha_1+\alpha_3)\alpha_2 = \alpha_1\alpha_2 \Leftrightarrow \alpha_2 = 0 \vee \alpha_3 = 0$, a contradiction. Therefore, $x\notin U^{i}U^{i+1}$.
    This proves that, for $q\neq 2$, $f=3$ is the minimal positive integer such that 
    $U^{[i,i+1]}_f\cap U^{[i+1,i]}_f\neq U^{[i,i+1]}_{f-1}\cup U^{[i+1,i]}_{f-1}$.
    
    \textit{Case $q=2$:} As we are forcing that all $\alpha_k\neq 0$, with $q=2$ there is only one option, i.e. $\alpha_k=1$ for all $k$. However, this means that $\alpha_5 = 1 + 1 = 0$, a contradiction. Indeed, at least one of ${\alpha_1,\alpha_2,\alpha_3}$ and at least one of ${\alpha_4,\alpha_5,\alpha_6}$ must be zero.
    Therefore, all elements in the intersection $U^iU^{i+1}U^i\cap U^{i+1}U^{i}U^{i+1}$ can be written in either $U^iU^{i+1}$ or $U^{i+1}U^{i}$.
    Therefore, contrary to the case $q\neq 2$, whenever $q=2$ we have $ U^iU^{i+1}U^i\cap U^{i+1}U^{i}U^{i+1}= U^iU^{i+1}\cup U^{i+1}U^i$, i.e. $f=3$ is \emph{not} the minimal positive integer such that $U^{[i,i+1]}_f\cap U^{[i+1,i]}_f\neq U^{[i,i+1]}_{f-1}\cup U^{[i+1,i]}_{f-1}$. 

    Finally, we will prove that, whenever $q=2$, $f=4$ is the minimal positive integer satisfying $U^{[i,i+1]}_f\cap U^{[i+1,i]}_f\neq U^{[i,i+1]}_{f-1}\cup U^{[i+1,i]}_{f-1}$. Take $x\in U^iU^{i+1}U^iU^{i+1}\cap U^{i+1}U^{i}U^{i+1}U^i$  such that $$x=(e_i(1)e_{i+1}(1))^2 = (e_{i+1}(1)e_i(1))^2=e_{i,i+2}(1).$$
    An element of $U^iU^{i+1}U^i$ is always written as $$e_i(\alpha_1)e_{i+1}(\alpha_2)e_{i}(\alpha_3)=\mathbf{1}+f_i(\alpha_1+\alpha_3)+f_{i+1}(\alpha_2)+f_{i,i+2}(\alpha_1\alpha_2).$$
    For $x$ to be in $U^iU^{i+1}U^i$, we would need $\alpha_1\alpha_2=1$ and $\alpha_2=0$, a contradiction. The same argument can be said for elements of $U^{i+1}U^iU^{i+1}$. Hence, $x=e_{i,i+2}(1)\notin U^iU^{i+1}U^i\cup U^{i+1}U^iU^{i+1}$. Therefore, for $q\neq 2$, $f=4$ is the minimal positive integer such that 
    $U^{[i,i+1]}_f\cap U^{[i+1,i]}_f\neq U^{[i,i+1]}_{f-1}\cup U^{[i+1,i]}_{f-1}$.
\end{proof}

\end{document}
